\numberwithin{equation}{section}

\newtheorem{Theorem}{Theorem}[section]
\newtheorem{Corollary}[Theorem]{Corollary}
\newtheorem{Lemma}[Theorem]{Lemma}
\newtheorem{Proposition}[Theorem]{Proposition}
{ \theoremstyle{definition}
\newtheorem{Definition}[Theorem]{Definition}
\newtheorem{Note}[Theorem]{Note}
\newtheorem{Example}[Theorem]{Example}
\newtheorem{Remark}[Theorem]{Remark}
\newtheorem{Notation}[Theorem]{Notation}
\newtheorem{Assumptions}[Theorem]{Assumptions}
}

\usepackage{mathtools}

%\usepackage{mathtools, MnSymbol, stmaryrd}
\usepackage{tikz}
\usetikzlibrary{decorations.pathreplacing, calc, shapes.geometric, shapes.misc, cd, positioning}
%\usepackage{IEEEtrantools}
%\usepackage{mleftright}
%\usepackage{csquotes}
\usepackage{tabularx}
%\usepackage[shortcuts]{extdash}

\usepackage[only,llbracket,rrbracket]{stmaryrd}

% Wrapper glyph compatibility: diamond with central plus, replacing unavailable MnSymbolC font.
\newcommand{\diamondplus}{\mathbin{\ooalign{\hfil$\diamond$\hfil\cr\hfil$\scriptscriptstyle+$\hfil\cr}}}

\usepackage{enumitem}
\setlist[enumerate]{nosep}
\setlist[enumerate,1]{label=$(\alph*)$}
\setlist[enumerate,2]{label=(\roman*)}
\setlist[enumerate,3]{label=(\arabic*)}
\setlist[itemize]{nosep}
\setlist[itemize,1]{label=$\square$}
\newtheorem*{theorem*}{Theorem}
\newcommand\Tstrut{\rule{0pt}{2.6ex}}
\newcommand{\identitymatrix}{I}
\newcommand{\freealg}[2]{#1\langle #2\rangle}
\newcommand{\univalg}[3]{#1\langle #2 \classpredicate #3\rangle}
\newcommand{\HSalgebraconstructdeform}[4]{#2\diamondplus^{#1}_{#4}#3}
\newcommand{\HScocycles}[3]{Z^{#1}_{\mathrm{HS}}(#2,#3)}
\newcommand{\HScoboundaries}[3]{B^{#1}_{\mathrm{HS}}(#2,#3)}
\newcommand{\HScohomology}[3]{H^{#1}_{\mathrm{HS}}(#2,#3)}
\newcommand{\HSdifferential}[1]{\partial^{#1}}
\newcommand{\HSuacocycles}[4]{Z^{#1}_{#2,#3,#4}}
\newcommand{\HSuacoboundaries}[4]{B^{#1}_{#2,#3,#4}}
\newcommand{\HSuacohomology}[4]{H^{#1}_{#2,#3,#4}}
\newcommand{\HSuafunc}[4]{F^{#1}_{#2,#3,#4}}
\newcommand{\HSuafuncx}[5]{F^{#1,#5}_{#2,#3,#4}}
\newcommand{\HSuafuncxx}[6]{F^{#1,#5}_{#2,#3,#4}(#6)}
\newcommand{\HSuafuncxxx}[7]{F^{#1,#5}_{#2,#3,#4}(#6)(#7)}
\newcommand{\thefunc}{F}
\newcommand{\thefuncx}[1]{F_{#1}}
\newcommand{\thefuncxx}[2]{F_{#1}(#2)}
\newcommand{\Scomultiplication}{\Delta}
\newcommand{\Scounit}{\epsilon}
\newcommand{\Sstar}{\ast}
\newcommand{\tiff}{\text{if and only if}}
\newcommand{\tif}{\text{if }}
\newcommand{\tothw}{\text{otherwise}}
\newcommand{\tand}{\text{and}}
\newcommand{\tor}{\text{or}}
\newcommand{\twhere}{\text{where}}
\newcommand{\twith}{\text{with}}
\newcommand{\todd}{\text{odd}}
\newcommand{\teven}{\text{even}}
\newcommand{\Sand}{\,\wedge\,}
\newcommand{\Sor}{\,\vee\,}
\newcommand{\Sthen}{\,\Rightarrow\,}
\newcommand{\Simplies}{\,\implies\,}
\newcommand{\Siff}{\Leftrightarrow}
\newcommand{\Strue}{\top}
\newcommand{\Sfalse}{\bot}
%\newcommand{\eqpd}{\coloneqq}
\newcommand{\iffpd}{{\ \colon\!\!\!\iff}}
\newcommand{\funcdef}{\colon}
\newcommand{\classpredicate}{\mid}
\newcommand{\quantorpredicate}{\colon }
\newcommand{\Sidentity}{\mathrm{id}}
\newcommand{\SidentityX}[1]{\Sidentity_{#1}}
\newcommand{\Scomposition}{\circ}
\newcommand{\Smonoidalproduct}{\otimes}
\newcommand{\Sinternalhom}{[\argph,\argph]}
\newcommand{\SinternalhomX}[2]{[#1,#2]}


\newcommand{\nnint}{{\mathbb{N}_0}}
\newcommand{\pint}{\mathbb{N}}
\newcommand{\integers}{\mathbb{Z}}
\newcommand{\comps}{\mathbb{C}}
\newcommand{\reals}{\mathbb{R}}
\newcommand{\dwi}[1]{\llbracket{#1}\rrbracket}
\newcommand{\fromto}[2]{#1 \to #2}
\newcommand{\xfromto}[3]{#1 \colon  #2 \to #3}
\newcommand{\fromtomaps}[4]{#1 \to #2, \, #3 \mapsto #4}
\newcommand{\xfromtomaps}[5]{#1 \colon  #2 \to #3, \, #4 \mapsto #5}
\newcommand{\sran}{\mathrm{ran}}
\newcommand{\sranx}[1]{\mathrm{ran}(#1)}
\newcommand{\bfunc}{,\,}

\newcommand{\ASdualX}[1]{\overline{#1}}
\newcommand{\subproof}[1]{\textit{#1}}
\newcommand{\Sisomorphic}{\cong}
\newcommand{\Sdirectsum}{\oplus}
\newcommand{\Sdirectproduct}{\times}
\newcommand{\Ssetmonoidalproduct}{\times}
\newcommand{\Sleftmoduleaction}{\triangleright}
\newcommand{\Srightmoduleaction}{\triangleleft}
\newcommand{\Sadj}{^\ast}
\newcommand{\Stra}{^{\mathrm{t}}}
\newcommand{\argph}{\hspace{0.5pt}\cdot\hspace{0.5pt}}

\newcommand{\upp}[1]{\prescript{\resizebox{0.008\hsize}{!}{$\blacksquare$}}{}{#1}}
\newcommand{\lop}[1]{\prescript{}{\resizebox{0.008\hsize}{!}{$\blacksquare$}}{#1}}
\newcommand{\tsopx}[2]{\Pi^{#1}_{#2}}
\newcommand{\blpoint}{\mathnormal{\bullet}}
\newcommand{\whpoint}{\mathnormal{\circ}}
\newcommand{\blaw}{\{\whpoint,\blpoint\}}
\newcommand{\tcsx}[2]{\Sigma^{#1}_{#2}}
\newcommand{\csx}[3]{\sigma^{#1}_{#2}(#3)}
\newcommand{\csfunx}[2]{\sigma^{#1}_{#2}}
\newcommand{\simg}[1]{#1_{\rightarrow}}
\newcommand{\simgx}[2]{#1_{\rightarrow}(#2)}
\newcommand{\spimg}[1]{#1^{\leftarrow}}
\newcommand{\spimgx}[2]{#1^{\leftarrow}(#2)}
\newcommand{\finerthan}{\mathrel{\leq}}
\newcommand{\nfinerthan}{\mathrel{\not\leq}}
\newcommand{\finerthanrel}{\mathop{\leq}}
\newcommand{\Yc}[1]{\mathfrak{#1}}
\newcommand{\Yp}[1]{\mathtt{#1}}
\newcommand{\Yb}[1]{\mathsf{#1}}
\newcommand{\blo}[2]{\pi_{#1}(#2)}
\newcommand{\blofun}[1]{\pi_{#1}}
\newcommand{\kerp}[1]{\ker(#1)}
\newcommand{\djp}[2]{#1\begin{smallmatrix}\resizebox{0.008\hsize}{!}{$\blacksquare$}\hspace{4pt}\\\hspace{4pt}\resizebox{0.008\hsize}{!}{$\blacksquare$}\end{smallmatrix}#2}
\newcommand{\zetf}{\zeta}
\newcommand{\zetfx}[2]{\zetf(#1,#2)}
\newcommand{\squarematrices}[2]{M_{#1}(#2)}
\newcommand{\kron}[2]{\delta_{#1,#2}}
\newcommand{\quomap}[2]{\sfrac{#1}{#2}}
\newcommand{\tsquomap}[2]{#1 / #2}
\newcommand{\onehalf}{{ \frac{1}{2}}}
\newcommand{\rightwardproduct}[2]{\underset{#1}{\overset{#2}{\overrightarrow{\prod}}}}
\newcommand{\tsrightwardproduct}[2]{{\overrightarrow{ \prod}}\vphantom{ \prod}_{#1}^{#2}}
\newcommand{\leftwardproduct}[2]{\underset{#1}{\overset{#2}{\overleftarrow{\prod}}}}
\newcommand{\tsleftwardproduct}[2]{{\overleftarrow{ \prod}}\vphantom{ \prod}_{#1}^{#2}}

\newcommand{\DQGgroupalgebra}[1]{\comps [#1 ]}
\newcommand{\DQGcocycles}[3]{Z^{#1} (#2,#3 )}
\newcommand{\DQGcoboundaries}[3]{B^{#1} (#2,#3 )}
\newcommand{\DQGcohomology}[3]{H^{#1} (#2,#3 )}
\newcommand{\DQGcocyclesTC}[2]{Z^{#1} (#2 )}
\newcommand{\DQGcoboundariesTC}[2]{B^{#1} (#2 )}
\newcommand{\DQGcohomologyTC}[2]{H^{#1} \bigl(#2 \bigr)}
\newcommand{\CQGdualC}[1]{\widehat{#1}}

\newcommand{\ueqgU}[1]{U_{#1}}
\newcommand{\ueqgUplus}[1]{U^{+}_{#1}}
\newcommand{\ueqgUstar}[2]{U^{\ast}_{#1,#2}}
\newcommand{\ueqgUtimes}[2]{U^{\times}_{#1,#2}}
\newcommand{\ueqgUtimesplus}[2]{U^{{\times}{+}}_{#1,#2}}



\newcommand{\Slength}[1]{|#1|}
\newcommand{\therelpoly}[5]{\mathrm{r}^{#1}_{#2}(#3)_{#4,#5}}

\newcommand{\alterer}[3]{#1\downarrow_{#2}#3}

\newcommand{\thedim}{n}
\newcommand{\theqg}{G}
\newcommand{\thegens}{E}
\newcommand{\theuni}{u}
\newcommand{\theunim}[1]{\theuni^{#1}}
\newcommand{\theunix}[3]{\theuni^{#1}_{#2,#3}}
\newcommand{\therels}{R}
\newcommand{\thepartrels}[1]{\therels_{#1}}
\newcommand{\thecat}{\mathcal{C}}
\newcommand{\theideal}{J}
\newcommand{\thepartideal}[1]{\theideal_{#1}}
\newcommand{\thefield}{\mathbb{K}}
\newcommand{\themodule}{X}
\newcommand{\theresmodule}{Y}
\newcommand{\thecounit}{\Scounit}
  \newcommand{\thepredicate}{A}
  \newcommand{\thepredicatex}[2]{A(#1,#2)}

  \newcommand{\scenarioone}{R1}
  \newcommand{\scenariotwo}{R2}
  \newcommand{\scenariothree}{R3}
  \newcommand{\scenariofour}{R4}

  \newcommand{\conditionone}{P1}
  \newcommand{\conditiontwo}{P2}
  \newcommand{\conditionthree}{P3}

\newcommand{\oeqgS}[1]{S_{#1}}
\newcommand{\oeqgSprime}[1]{S^{\prime}_{#1}}
\newcommand{\oeqgSplus}[1]{S^{+}_{#1}}
\newcommand{\oeqgSprimeplus}[1]{S^{\prime +}_{#1}}

\newcommand{\oeqgO}[1]{O_{#1}}
\newcommand{\oeqgOstar}[1]{O^{\ast}_{#1}}
\newcommand{\oeqgOplus}[1]{O^{+}_{#1}}

\newcommand{\oeqgB}[1]{B_{#1}}
\newcommand{\oeqgBprime}[1]{B^{\prime}_{#1}}
\newcommand{\oeqgBhashstar}[1]{B^{\#\ast}_ {#1}}
\newcommand{\oeqgBplus}[1]{B^{+}_{#1}}
\newcommand{\oeqgBprimeplus}[1]{B^{\prime{+}}_{#1}}
\newcommand{\oeqgBhashplus}[1]{B^{\#{+}}_{#1}}

\newcommand{\oeqgH}[1]{H_{#1}}
\newcommand{\oeqgHstar}[1]{H^{\ast}_{#1}}
\newcommand{\oeqgHplus}[1]{H^{+}_{#1}}
\newcommand{\oeqgHcurly}[2]{H^{\{#1\}}_{#2}}
\newcommand{\oeqgHcurlyinfinity}[1]{H^{\{\infty\}}_{#1}}
\newcommand{\oeqgHsquare}[2]{H^{[#1]}_{#2}}
\newcommand{\oeqgHsquareinfinity}[1]{H^{[\infty]}_{#1}}
\newcommand{\oeqgHround}[2]{H^{(#1)}_{#2}}
\newcommand{\oeqgHangle}[2]{H^{\langle #1\rangle}_{#2}}




  \newcommand{\onlyneutralnonsingletonblocks}{\mathsf{NNSB}}
  \newcommand{\onlyneutralpartitions}{\mathsf{NP}}

\newcommand{\PartIdenB}{%
  \begin{tikzpicture}[scale=0.25,baseline=0.015cm]
    \def\xdist{1}
    \def\ydist{1}
    \node [scale=0.33, circle, draw=black, fill=black] (a1) at ({0*\xdist},{0*\ydist}) {};
    \node [scale=0.33, circle, draw=black, fill=black] (b1) at ({0*\xdist},{1*\ydist}) {};
    \draw (a1) to (b1);
  \end{tikzpicture}
}
\newcommand{\PartIdenW}{%
  \begin{tikzpicture}[scale=0.25,baseline=0.015cm]
    \def\xdist{1}
    \def\ydist{1}
    \node [scale=0.33, circle, draw=black, fill=white] (a1) at ({0*\xdist},{0*\ydist}) {};
    \node [scale=0.33, circle, draw=black, fill=white] (b1) at ({0*\xdist},{1*\ydist}) {};
    \draw (a1) to (b1);
  \end{tikzpicture}
}
\newcommand{\PartIdenUpBW}{%
  \begin{tikzpicture}[scale=0.25,baseline=-0.015cm]
    \def\xdist{1}
    \def\ydist{1}
    \def\ypar{1}
    \node [scale=0.33, circle, draw=black, fill=black] (a1) at ({0*\xdist},{1*\ydist}) {};
    \node [scale=0.33, circle, draw=black, fill=white] (a2) at ({1*\xdist},{1*\ydist}) {};
    \draw (a1) -- ++ (0,{-\ypar}) -| (a2);
  \end{tikzpicture}
}
\newcommand{\PartIdenUpWB}{%
  \begin{tikzpicture}[scale=0.25,baseline=-0.015cm]
    \def\xdist{1}
    \def\ydist{1}
    \def\ypar{1}
    \node [scale=0.33, circle, draw=black, fill=white] (a1) at ({0*\xdist},{1*\ydist}) {};
    \node [scale=0.33, circle, draw=black, fill=black] (a2) at ({1*\xdist},{1*\ydist}) {};
    \draw (a1) -- ++ (0,{-\ypar}) -| (a2);
  \end{tikzpicture}
  }
\newcommand{\PartIdenLoBW}{%
  \begin{tikzpicture}[scale=0.25,baseline=-0.015cm]
    \def\xdist{1}
    \def\ydist{1}
    \def\ypar{1}
    \node [scale=0.33, circle, draw=black, fill=black] (a1) at ({0*\xdist},{0*\ydist}) {};
    \node [scale=0.33, circle, draw=black, fill=white] (a2) at ({1*\xdist},{0*\ydist}) {};
    \draw (a1) -- ++ (0,{\ypar}) -| (a2);
  \end{tikzpicture}
}
\newcommand{\PartIdenLoWB}{%
  \begin{tikzpicture}[scale=0.25,baseline=-0.015cm]
    \def\xdist{1}
    \def\ydist{1}
    \def\ypar{1}
    \node [scale=0.33, circle, draw=black, fill=white] (a1) at ({0*\xdist},{0*\ydist}) {};
    \node [scale=0.33, circle, draw=black, fill=black] (a2) at ({1*\xdist},{0*\ydist}) {};
    \draw (a1) -- ++ (0,{\ypar}) -| (a2);
  \end{tikzpicture}
}
\newcommand{\PartIdenLoWW}{%
  \begin{tikzpicture}[scale=0.25,baseline=-0.015cm]
    \def\xdist{1}
    \def\ydist{1}
    \def\ypar{1}
    \node [scale=0.33, circle, draw=black, fill=white] (a1) at ({0*\xdist},{0*\ydist}) {};
    \node [scale=0.33, circle, draw=black, fill=white] (a2) at ({1*\xdist},{0*\ydist}) {};
    \draw (a1) -- ++ (0,{\ypar}) -| (a2);
  \end{tikzpicture}
}
\newcommand{\PartFourWBWB}{%
  \begin{tikzpicture}[scale=0.25,baseline=-0.015cm]
    \def\xdist{0.666}
    \def\ydist{1}
    \def\ypar{1}
    \node [scale=0.33, circle, draw=black, fill=white] (a1) at ({0*\xdist},{0*\ydist}) {};
    \node [scale=0.33, circle, draw=black, fill=black] (a2) at ({1*\xdist},{0*\ydist}) {};
    \node [scale=0.33, circle, draw=black, fill=white] (a3) at ({2*\xdist},{0*\ydist}) {};
    \node [scale=0.33, circle, draw=black, fill=black] (a4) at ({3*\xdist},{0*\ydist}) {};
    \draw (a1) -- ++ (0,{\ypar}) -| (a4);
    \draw (a2) -- ++ (0,{\ypar});
    \draw (a3) -- ++ (0,{\ypar});
  \end{tikzpicture}
}
\newcommand{\PartSinglesWB}{%
  \begin{tikzpicture}[scale=0.25,baseline=-0.015cm]
    \def\xdist{1}
    \def\ydist{1}
    \def\ypar{1}
    \node [scale=0.33, circle, draw=black, fill=white] (a1) at ({0*\xdist},{0*\ydist}) {};
    \node [scale=0.33, circle, draw=black, fill=black] (a2) at ({1*\xdist},{0*\ydist}) {};
    \draw[->] (a1) -- ++ (0,{\ypar});
    \draw[->] (a2) -- ++ (0,{\ypar});
  \end{tikzpicture}
}
\newcommand{\PartSingleW}{%
    \begin{tikzpicture}[scale=0.25,baseline=-0.015cm]
    \def\xdist{1}
    \def\ydist{1}
    \def\ypar{1}
    \node [scale=0.33, circle, draw=black, fill=white] (a1) at ({0*\xdist},{0*\ydist}) {};
    \draw[->] (a1) -- ++ (0,{\ypar});
  \end{tikzpicture}
}
\newcommand{\PartSingleB}{%
    \begin{tikzpicture}[scale=0.25,baseline=-0.015cm]
    \def\xdist{1}
    \def\ydist{1}
    \def\ypar{1}
    \node [scale=0.33, circle, draw=black, fill=black] (a1) at ({0*\xdist},{0*\ydist}) {};
    \draw[->] (a1) -- ++ (0,{\ypar});
  \end{tikzpicture}
}
\newcommand{\PartCoSingleWBWB}{%
    \begin{tikzpicture}[scale=0.25,baseline=-0.015cm]
    \def\xdist{0.666}
    \def\ydist{1.25}
    \def\ypar{1}
    \node [scale=0.33, circle, draw=black, fill=white] (a1) at ({0*\xdist},{0*\ydist}) {};
    \node [scale=0.33, circle, draw=black, fill=black] (a2) at ({1*\xdist},{0*\ydist}) {};
    \node [scale=0.33, circle, draw=black, fill=white] (a3) at ({2*\xdist},{0*\ydist}) {};
    \node [scale=0.33, circle, draw=black, fill=black] (a4) at ({3*\xdist},{0*\ydist}) {};
    \draw (a1) -- ++ (0,{1*\ydist}) -| (a3);
    \draw[->] (a2) -- ++ (0,{1*\ypar});
    \draw[->] (a4) -- ++ (0,{1*\ypar});
    \end{tikzpicture}
  }
\newcommand{\PartCrossWW}{%
  \begin{tikzpicture}[scale=0.25,baseline=0.015cm]
    \def\xdist{1}
    \def\ydist{1}
    \node [scale=0.33, circle, draw=black, fill=white] (a1) at ({0*\xdist},{0*\ydist}) {};
    \node [scale=0.33, circle, draw=black, fill=white] (a2) at ({1*\xdist},{0*\ydist}) {};
    \node [scale=0.33, circle, draw=black, fill=white] (b1) at ({0*\xdist},{1*\ydist}) {};
    \node [scale=0.33, circle, draw=black, fill=white] (b2) at ({1*\xdist},{1*\ydist}) {};
    \draw (a1) to (b2);
    \draw (a2) to (b1);
  \end{tikzpicture}
}
\newcommand{\UCPartIden}{%
  \begin{tikzpicture}[baseline=0.0cm]
    \def\xdist{0.2cm}
    \def\ydist{0.32cm}
    \def\ypar{0.25cm}
    \coordinate (a1) at ({0*\xdist},{0*\ydist});
    \draw (a1) -- ++ (0,{\ypar});
  \end{tikzpicture}
}
\newcommand{\UCPartIdenLo}{%
  \begin{tikzpicture}[scale=0.25,baseline=0.015cm]
    \def\xdist{1}
    \def\ydist{1}
    \def\ypar{1}
    \coordinate (a1) at ({0*\xdist},{0*\ydist}) {};
    \coordinate (a2) at ({1*\xdist},{0*\ydist}) {};
    \draw (a1) -- ++ (0,{\ypar}) -| (a2);
  \end{tikzpicture}
  }
\newcommand{\UCPartIdenUp}{%
  \begin{tikzpicture}[scale=0.25,baseline=0.015cm]
    \def\xdist{1}
    \def\ydist{1}
    \def\ypar{1}
    \coordinate (a1) at ({0*\xdist},{1*\ydist}) {};
    \coordinate (a2) at ({1*\xdist},{1*\ydist}) {};
    \draw (a1) -- ++ (0,{-\ypar}) -| (a2);
  \end{tikzpicture}
  }
\newcommand{\UCPartFour}{%
  \begin{tikzpicture}[baseline=0.0cm]
    \def\xdist{0.2cm}
    \def\ydist{0.32cm}
    \def\ypar{0.2cm}
    \coordinate (a1) at ({0*\xdist},{0*\ydist});
    \coordinate (a2) at ({1*\xdist},{0*\ydist});
    \coordinate (a3) at ({2*\xdist},{0*\ydist});
    \coordinate (a4) at ({3*\xdist},{0*\ydist});
    \draw (a1) -- ++ (0,{\ypar}) -| (a4);
    \draw (a2) -- ++ (0,{\ypar});
    \draw (a3) -- ++ (0,{\ypar});
  \end{tikzpicture}
}
\newcommand{\UCPartCoSingle}{%
  \begin{tikzpicture}[baseline=0.0cm]
    \def\xdist{0.2cm}
    \def\ydist{0.32cm}
    \def\ypar{0.2cm}
    \coordinate (a1) at ({0*\xdist},{0*\ydist});
    \coordinate (a2) at ({1*\xdist},{0*\ydist});
    \coordinate (a3) at ({2*\xdist},{0*\ydist});
    \coordinate (a4) at ({3*\xdist},{0*\ydist});
    \draw (a1) -- ++ (0,{\ydist}) -| (a3);
    \draw[->] (a2) -- ++ (0,{\ypar});
    \draw[->] (a4) -- ++ (0,{\ypar});
  \end{tikzpicture}
}
\newcommand{\UCPartSingle}{%
  \begin{tikzpicture}[baseline=0.0cm]
    \def\xdist{0.2cm}
    \def\ydist{0.32cm}
    \def\ypar{0.25cm}
    \coordinate (a1) at ({0*\xdist},{0*\ydist});
    \draw[->] (a1) -- ++ (0,{\ypar});
  \end{tikzpicture}
}
\newcommand{\UCPartSinglesTensor}{%
  \begin{tikzpicture}[baseline=0.0cm]
    \def\xdist{0.2cm}
    \def\ydist{0.32cm}
    \def\ypar{0.25cm}
    \coordinate (a1) at ({0*\xdist},{0*\ydist});
    \coordinate (a2) at ({1.75*\xdist},{0*\ydist});
    \draw[->] (a1) -- ++ (0,{\ypar});
    \draw[->] (a2) -- ++ (0,{\ypar});
    \node [scale=0.66] at ($(a1)+({0.5*1.75*\xdist},{0.325*\ydist})$) {$\otimes$};
  \end{tikzpicture}
}
\newcommand{\UCPartSingles}{%
  \begin{tikzpicture}[baseline=0.0cm]
    \def\xdist{0.2cm}
    \def\ydist{0.32cm}
    \def\ypar{0.2cm}
    \coordinate (a1) at ({0*\xdist},{0*\ydist});
    \coordinate (a2) at ({1*\xdist},{0*\ydist});
    \draw[->] (a1) -- ++ (0,{\ypar});
    \draw[->] (a2) -- ++ (0,{\ypar});
  \end{tikzpicture}
}



% Local text-accent fallback: DejaVu Serif includes combining cedilla.
\usepackage{fontspec}
\newfontfamily\OAaccentfont{DejaVu Serif}
\let\OAoriginalcedilla\c
\renewcommand{\c}[1]{{\OAaccentfont\OAoriginalcedilla{#1}}}
